\begin{center}\begin{tikzpicture}[font=\small,>=stealth,line join=round]
\draw[gray!40,rounded corners=3pt] (0,0) rectangle (14.2,9.0);
\node at (7.1,8.4) {反证：右侧足够近的点，比左段末点更低};
\begin{scope}[shift={(3,3)},x=0.95cm,y=0.95cm]\draw[->,gray] (-2,0)--(2.3,0) node[right]{$x$};\draw[->,gray] (0,-2)--(0,4.3) node[above]{$y$};\begin{scope}\clip (-2,-2) rectangle (2,4);\draw[blue,thick,domain=-2:0,samples=80] plot(\x,{\x+2});\draw[orange!85!black,thick,domain=0.0001:2,samples=100] plot(\x,{\x});\end{scope}\draw[draw=blue,fill=blue] (0,2) circle(2pt);\draw[draw=orange!85!black,fill=white] (0,0) circle(2pt);\end{scope}\begin{scope}[shift={(3,3)},x=.95cm,y=.95cm]\fill[red!75!black] (.5,.5) circle(2.5pt) node[right]{$(1/2,1/2)$};\draw[red!70!black,dashed,->] (0,2)--(.5,.5);\end{scope}\node[align=left,anchor=west] at (7.3,4.1) {$f(0)=2$，$L=0$\\[12pt]取 $x_2=1/2>0$\\[10pt]却有 $f(x_2)=1/2<2$\\[12pt]一对逆序点即可否定整体递增。};
\node[text width=12.8cm,align=center] at (7.1,.65) {右段每点都在上升，但跨过分界点仍能找到“向右却降低”的一对点。};
\end{tikzpicture}\end{center}
